\chapter{Conclusion}
\label{ch:conclusion}

\section{Summary of Contributions}

This thesis set out to investigate whether multi-agent LLM architectures could provide intrusion detection that is simultaneously accurate enough to be operationally useful and transparent enough to address the explainability deficit of traditional machine learning classifiers. The experimental results provide an affirmative answer to both halves of that question.

The first contribution is the first systematic, per-attack-type evaluation of a multi-agent LLM intrusion detection system at realistic class distributions. Prior evaluations of LLM-based NIDS used balanced datasets or aggregated multi-attack batches, neither of which reflects the class imbalance of production networks. The Stage 1 evaluation reveals a performance range from F1=100\% on brute-force and DoS categories to F1=0\% on Infiltration --- a range that would be invisible in aggregate evaluation.

The second contribution is a two-tier ML-plus-LLM architecture reducing the cost of LLM-based detection by 94.6 per cent. Actual expenditure across all fourteen Stage 1 experiments was \$27.35, against an estimated \$509.78 without Tier 1 pre-filtering. The Random Forest at threshold 0.15 achieves one hundred per cent recall on attack flows in the development partition, ensuring that cost reduction is not purchased at the expense of missed detections. This two-tier pattern is generalisable to any domain where LLM analysis is accurate but expensive and a cheaper discriminator can identify the easy cases.

The third contribution is the introduction of a Devil's Advocate agent as a structural false positive control mechanism, demonstrated to achieve zero per cent FPR across thirteen of fourteen attack categories. The calibration of the Devil's Advocate weight at thirty per cent provides a concrete, reusable design parameter for practitioners building adversarial multi-agent security systems, extending the empirical grounding of adversarial agent design \autocite{Yin2024, Du2023, Chan2023} into quantified network security performance.

The fourth contribution is a complete explainability framework where every detection decision is accompanied by multi-perspective reasoning chains from six analytical agents, operationalised through the Flow Inspector interface. Unlike SHAP or LIME post-hoc explanations, AMATAS reasoning chains are generated as part of the detection process itself, grounded in domain knowledge, and structured to present both evidence for malice and the best available counter-argument.

The fifth contribution is a controlled ablation study quantifying the marginal value of each architectural addition --- from zero-shot single agent (F1=62.8\%) to prompt-engineered single agent (F1=58.0\%) to tool-augmented single agent (F1=58.3\%) to multi-agent AMATAS (F1=88\%) --- and confirming that external threat intelligence tools provide no measurable benefit on anonymised datasets.

\section{Answers to the Research Questions}

\textbf{RQ1} asked whether a multi-agent LLM architecture can achieve detection accuracy comparable to traditional machine learning classifiers. AMATAS achieves mean F1 of 92.9 per cent across thirteen detected attack types, with four categories at F1=100\%. This does not match the theoretical ceiling of fine-tuned discriminative models but is operationally competitive with deployed traditional classifiers, which degrade under distribution shift and novel attack classes. More importantly, AMATAS achieves this performance whilst providing explanations that discriminative classifiers cannot; on that framing, RQ1 is answered affirmatively.

\textbf{RQ2} asked how a two-tier architecture affects the cost-accuracy trade-off. The Tier 1 Random Forest filters approximately ninety-four per cent of flows as confidently benign, reducing Stage 1 cost from an estimated \$509.78 to an actual \$27.35 --- a 94.6 per cent reduction --- whilst maintaining one hundred per cent recall on attack flows at the operating threshold. The resulting cost of approximately \$0.002 per flow makes LLM-based analysis economically viable for batch processing of millions of flows.

\textbf{RQ3} asked about the quality and utility of multi-agent reasoning chains relative to opaque classifier outputs. The reasoning chains generated by AMATAS are qualitatively different from feature attribution outputs in several respects: they are generated as part of detection rather than as post-hoc approximations; they incorporate domain knowledge about attack signatures, protocol semantics, and MITRE ATT\&CK techniques; they present structured disagreement between analytical perspectives; and they include an explicit counter-argument. These structural properties --- specificity, grounding, multi-perspectival structure, and falsifiability --- support a positive answer to RQ3.

\textbf{RQ4} asked how detection performance varies across attack types and what the fundamental limitations of flow-level LLM analysis are. Performance ranges from F1=100\% for attacks with distinctive flow signatures to F1=0\% for Infiltration, where individual attack flows are statistically indistinguishable from benign traffic. No single-flow analysis method --- LLM or otherwise --- can reliably detect attacks that have been designed to resemble benign traffic at the individual flow level.

\section{Future Work}
\label{sec:future}

The most immediately motivated extension is AMATAS v3, the temporal clustering architecture designed to address the Infiltration failure and improve recall on DDoS-HOIC. The central hypothesis is that grouping temporally adjacent flows from the same source IP into ordered analysis units will reveal the access progression patterns that characterise infiltration attacks. Implementation requires extending the batch creation pipeline to produce sequence-aware input structures, modifying the Temporal Agent to reason over ordered sequences rather than unordered co-IP flow sets, and re-evaluating specifically on Infiltration and DDoS-HOIC categories.

A systematic agent ablation study would substantially strengthen the architectural argument. The current evaluation compares multi-agent AMATAS against single-agent baselines but does not quantify the marginal contribution of each additional agent --- the gain from the Temporal Agent, the gain from the Devil's Advocate, and the gain from the Orchestrator's structured synthesis relative to majority voting. Running two-agent through six-agent configurations on the same batch would precisely localise where the multi-agent benefit is concentrated.

Evaluation across alternative frontier LLMs is a priority for external validity. Claude Sonnet 4, Llama-3.3-70B, and Gemini 1.5 Pro are candidates for comparative evaluation; each would require re-calibration of the Devil's Advocate weight and Orchestrator consensus thresholds against the model's specific reasoning and confidence patterns. \textcite{Ferrag2020} document the sensitivity of LLM-based detection systems to the specific model used, and this sensitivity has only increased as model behaviours have diverged.

A production deployment integrating MCP-based tool access against live traffic with real IP addresses represents a natural extension of the comparison experiments. The null result observed on anonymised data does not preclude substantial tool utility on production traffic, where AbuseIPDB reputation scores, geolocation data, and MITRE ATT\&CK technique lookups could provide contextual enrichment that the agents' parametric knowledge cannot. Such a deployment would require addressing the security considerations discussed in Section~\ref{sec:mcp-security}, particularly tool response verification and permission boundaries, but the MCP server infrastructure developed for the comparison experiments provides a ready foundation for this extension.

The held-out test partition --- 8.05 million flows excluded from all development and validation decisions --- represents the definitive evaluation of AMATAS generalisability and should be treated as a one-shot evaluation with no post-hoc parameter adjustment. Longer-term directions include LoRA-adapted specialist models replacing each agent with a fine-tuned variant trained on agent-specific analytical labels, a streaming deployment architecture for real-time flow processing with asynchronous LLM queuing, and a formal user study providing empirical grounding for the explainability claims that this thesis argues for on structural grounds.
